% =========================================================================
% paper_body.tex — shared content, included by both paper.tex (IEEEtran)
% and paper_preview.tex (local two-column preview).
% All [TBD: ...] markers are placeholders to be filled from metrics.json /
% benchmark_report.json after running the pipeline on the real datasets.
% =========================================================================

\begin{abstract}
Modern industrial systems generate massive volumes of continuous multivariate
time-series data from hundreds of sensors simultaneously, in which latent
equipment defects are masked by overlapping sensor patterns. Existing
analytical frameworks struggle to reconcile two conflicting engineering
requirements: processing high-velocity periodic streams with minimal
computational overhead, and achieving low-latency anomaly alerting. This
paper presents a scalable pipeline architecture for real-time anomaly
detection in industrial multivariate time series under a strict sub-20\,ms
service-level agreement (SLA). The architecture couples a multi-threaded,
back-pressure-aware streaming ingestion layer with two interchangeable
analytical backends: a reconstruction-based LSTM autoencoder, and a
fast-initializing mathematical backend combining seasonal-trend decomposition
(STL) with a closed-form Dynamic Mode Decomposition (DMD) operator. Both
backends feed a common scoring node that fuses localized
reconstruction/prediction error with the harmonic mean of Euclidean distance
to reference operating-regime centroids. We evaluate both backends on the
NASA C-MAPSS turbofan degradation benchmark and the SWaT secure water
treatment dataset, measuring detection accuracy (F1, AUCROC) and per-vector
scoring latency against the 20\,ms SLA under both steady-state and burst
throughput conditions. [TBD: one to two sentences summarizing the headline
empirical findings once Phase~4 benchmarks have been run.]
\end{abstract}

\begin{IEEEkeywords}
anomaly detection, multivariate time series, industrial IoT, streaming
architecture, autoencoders, dynamic mode decomposition, AIOps, low-latency
systems
\end{IEEEkeywords}

% =========================================================================
\section{Introduction}
% =========================================================================
\IEEEPARstart{M}{odern} industrial systems are highly interconnected,
generating continuous multivariate time-series (MTS) streams from hundreds of
sensors simultaneously. Operational complexity frequently masks latent
equipment defects, leading to unexpected system breakdowns and expensive
operational downtime. Detecting anomalous behavior in these streams quickly
enough to act on it is a distinct engineering problem from detecting it
accurately: a model that is highly accurate but too slow to keep pace with
the incoming stream is not deployable in a real-time industrial control loop.

Current analytical frameworks fail to reconcile two conflicting engineering
needs: processing massive periodic streams with minimal computational
overhead, and achieving low-latency alerting~\cite{tao2025realtime,
exarchos2021anomaly}. Deep reconstruction-based models such as LSTM
autoencoders offer strong detection accuracy on complex multi-sensor
data~\cite{malhotra2015lstm}, but impose per-inference compute costs and
prolonged initialization periods that are difficult to sustain in
fast-paced production ecosystems. Conversely, lightweight statistical
methods initialize quickly and score cheaply, but have historically been
overwhelmed by the volume, velocity, and periodicity of industrial streams.

This paper addresses that tension directly. We architect, implement, and
evaluate an end-to-end pipeline for anomaly detection in industrial MTS
data under a strict sub-20\,ms SLA, defined as the elapsed time from a
sensor vector leaving the ingestion queue to an anomaly score being
produced. The contributions are threefold:

\begin{enumerate}
  \item \textbf{An end-to-end, reproducible pipeline} spanning offline model
        development, multi-threaded streaming ingestion with explicit
        back-pressure semantics, and SLA-aware per-vector latency
        instrumentation.
  \item \textbf{A direct empirical comparison} of a deep reconstruction-based
        backend (LSTM autoencoder) against a closed-form mathematical
        backend (STL~+~DMD) on identical data, identical streaming
        infrastructure, and identical evaluation criteria --- isolating the
        accuracy/latency trade-off from confounding pipeline differences.
  \item \textbf{An elasticity analysis} of both backends under sudden
        throughput spikes, measuring drop rate and tail latency beyond
        steady-state benchmarking alone.
\end{enumerate}

The remainder of this paper is organized as follows. Section~II reviews
related work. Section~III describes the system architecture.
Section~IV details the experimental methodology. Section~V presents
results, and Section~VI concludes.

% =========================================================================
\section{Related Work}
% =========================================================================

\subsection{AIOps and Industrial MTS Anomaly Detection}
The reliance of modern industrial plants on automated machinery has made MTS
anomaly detection vital for infrastructure reliability. In dense industrial
settings, anomalies are hidden within overlapping sensor patterns, such as
spatial correlations between vibration, pressure, and temperature metrics.
Traditional statistical methods are frequently overwhelmed by the volume,
velocity, and periodic nature of these streams, producing delayed alerts or
elevated false-alarm rates~\cite{tao2025realtime, exarchos2021anomaly}.
Entropy-based sequence characterization~\cite{richman2000physiological}
established early foundations for detecting irregularity in physiological
and mechanical time series, but does not by itself address the throughput
requirements of modern industrial deployments.

\subsection{Reconstruction-Based Deep Learning}
Reconstruction-based neural models detect anomalies by learning a
compressed latent representation of normal operational cycles and flagging
inputs that reconstruct poorly. LSTM-based encoder--decoder architectures
extend this to sequential data by modeling temporal dependencies within
each window~\cite{malhotra2015lstm}, and have proven effective in complex
multi-sensor industrial fields. However, deep models present structural
challenges in production ecosystems: prolonged training and initialization
periods, and per-inference costs that scale poorly against tight latency
budgets on commodity hardware~\cite{exarchos2021anomaly}.

\subsection{Fast-Initializing Mathematical Alternatives}
A complementary line of work favors closed-form mathematical methods over
gradient-trained models, prioritizing rapid initialization and low,
predictable inference latency. Seasonal-trend decomposition via
Loess (STL)~\cite{cleveland1990stl} separates a series into trend,
seasonal, and residual components, allowing anomalies to surface as
elevated residual energy relative to a learned baseline. Dynamic Mode
Decomposition (DMD)~\cite{schmid2010dmd} fits a best-fit linear operator
approximating one-step-ahead dynamics through a single truncated singular
value decomposition, requiring no iterative training. Recent systems
combining these primitives --- exemplified by the VersaGuardian
model\footnote{[TBD: the VersaGuardian model (2026) is described in the
project proposal's background research with reported figures of
approximately 20-minute initialization and 15.28\,ms average detection
latency; a citable primary source must be located and inserted here before
submission. These figures should not be treated as established fact until
independently verified.]} --- report competitive detection accuracy at
initialization and latency costs far below deep alternatives.

\subsection{Research Gap}
Despite these methodological innovations, a persistent gap remains in
coupling fast detection algorithms with full-scale streaming data
architectures. An algorithm is only as useful as the pipeline supplying its
data: industrial platforms require structured streaming ingestion that
feeds high-throughput streams into an evaluation engine without introducing
serialization or I/O bottlenecks. Direct empirical comparisons of deep and
closed-form backends under an explicit real-time SLA --- accounting for
queueing mechanics, back-pressure, and throughput-spike elasticity rather
than offline batch accuracy alone --- remain comparatively rare. This work
addresses that gap.

% =========================================================================
\section{System Architecture}
% =========================================================================
The pipeline is organized into three programmatic stages: an ingestion
layer, a processing and analytics core, and an evaluation matrix
(scoring node). Fig.~\ref{fig:architecture} summarizes the data flow.

\begin{figure}[!t]
\centering
\fbox{\parbox{0.9\columnwidth}{\centering\vspace{2em}
[TBD: architecture diagram --- sensor window $\rightarrow$ parallel
Autoencoder / fast STL+DMD branches $\rightarrow$ shared scoring node
$\rightarrow$ anomaly decision. Export from project documentation.]
\vspace{2em}}}
\caption{End-to-end pipeline architecture. A single sensor window is scored
independently by each backend; both feed the same scoring node, isolating
the backend comparison from pipeline confounds.}
\label{fig:architecture}
\end{figure}

\subsection{Ingestion Layer}
The ingestion layer is a multi-threaded messaging-queue simulation in
Python. A producer thread reads pre-scaled sensor arrays and pushes
multivariate windows into a bounded, thread-safe queue at a configurable
production rate, simulating an industrial message bus. The bounded queue
makes producer back-pressure an explicit, observable property of the
system. Two overflow policies are supported: blocking (back-pressure
propagates to the producer) for steady-state operation, and
drop-on-full for elasticity testing, simulating a broker's overflow
behavior under a sudden throughput spike.

\subsection{Processing and Analytics Core}
A consumer thread pulls vectors from the queue and scores them one at a
time --- deliberately unbatched, so that measured latency reflects genuine
single-vector inference cost rather than amortized batch throughput. Two
interchangeable analytical backends are evaluated:

\subsubsection{LSTM Autoencoder}
A reconstruction-based encoder--decoder network trained exclusively on
normal-operation windows with early stopping on a held-out normal
validation split. The encoder compresses each window through stacked LSTM
layers into a latent vector; the decoder reconstructs the window from the
latent seed. At inference, the mean squared reconstruction error over the
window serves as the raw anomaly signal.

\subsubsection{Fast Mathematical Backend (STL~+~DMD)}
The fast backend combines two closed-form components fit once on normal
data: (i)~a per-channel STL residual-energy baseline, capturing the
expected trend/seasonal signature of healthy operation; and (ii)~a DMD
linear operator $A$ such that $x_{t+1} \approx A x_t$, obtained via a
single rank-truncated SVD over stacked snapshot pairs. At inference, the
operator's one-step-ahead prediction error, normalized by its baseline
error distribution on normal data, serves as the raw anomaly signal.
Neither component requires gradient-based training, yielding near-instant
initialization relative to the autoencoder.

\subsection{Evaluation Matrix (Scoring Node)}
Both backends' raw error signals are $z$-normalized against normal-data
baseline statistics and linearly combined with a multivariate distance
term: the harmonic mean of Euclidean distances between each in-window
reading and a set of $k$ reference ``normal operating regime'' centroids
obtained by $k$-means over normal training vectors. The harmonic mean is
used deliberately rather than the arithmetic mean: because it is dominated
by the smallest distances, a window's distance score remains low only if
its readings sit close to every relevant reference regime, so a single
nearby centroid cannot mask a genuinely anomalous vector. The combined
score for window $w$ with raw backend error $e(w)$ is
\begin{equation}
s(w) = \alpha \, \tilde{e}(w) + (1-\alpha)\, \widetilde{d_{\mathrm{hm}}}(w),
\label{eq:score}
\end{equation}
where $\tilde{\cdot}$ denotes $z$-normalization against normal-data
statistics and $d_{\mathrm{hm}}$ is the harmonic-mean distance. The
detection threshold is calibrated as a high percentile of the normal
validation split's own score distribution, requiring no anomaly labels.

% =========================================================================
\section{Experimental Methodology}
% =========================================================================

\subsection{Datasets}
Two publicly available industrial benchmarks are used.
\textbf{NASA C-MAPSS}~\cite{saxena2008damage} provides multivariate sensor
telemetry across the run-to-failure lifetime of simulated turbofan
engines; we recast it for anomaly detection by labeling the final $N$
cycles before failure of each test unit as anomalous, with healthy
training trajectories serving as the normal baseline.
\textbf{SWaT}~\cite{goh2016swat} is collected from a fully operational,
scaled-down secure water treatment testbed and provides natively labeled
process-control attacks; its attack file serves as the labeled test set
and its normal-operation file as the training baseline.

\subsection{Preprocessing}
Missing values are handled by forward/backward fill with a median
fallback. Continuous sensor metrics are normalized with MinMax scaling fit
\emph{exclusively} on normal-operation data, preventing anomalous-region
statistics from leaking into normalization. Cleaned readings are
structured into overlapping contiguous windows, the atomic unit scored by
both backends. A window inherits the anomalous label if any reading within
it is anomalous --- a conservative policy favoring recall on
partial-window anomalies.

\subsection{SLA Instrumentation}
A high-precision monotonic timer wraps exactly the span from a vector
leaving the ingestion queue to its anomaly score being produced, per
vector. Queue waiting time is deliberately excluded: it is a property of
the ingestion/back-pressure subsystem, whereas the 20\,ms SLA targets the
scoring engine's own computational cost. We report mean, median, 95th-,
and 99th-percentile latency, and the fraction of vectors violating the
SLA.

\subsection{Steady-State and Stress Protocols}
In the steady-state protocol, the producer emits vectors at a fixed target
rate through a large bounded queue with blocking overflow; both backends
process the identical stream. In the stress protocol, the producer emits
unthrottled bursts into a deliberately small queue with drop-on-full
overflow, and we report effective throughput, drop rate, queue high
watermark, and tail latency --- quantifying each backend's elasticity
under load beyond its steady-state envelope.

\subsection{Metrics}
Detection accuracy is reported as F1-score and area under the receiver
operating characteristic curve (AUCROC) at window level. Computational
efficiency is reported as the per-vector latency statistics above, tracked
against the 20\,ms SLA threshold.

% =========================================================================
\section{Results and Discussion}
% =========================================================================

\subsection{Detection Accuracy and Steady-State Latency}
Table~\ref{tab:steady} reports accuracy and latency for both backends on
both datasets under the steady-state protocol.
[TBD: populate from \texttt{metrics.json} and
\texttt{benchmark\_report.json} after running \texttt{train\_offline.py}
and \texttt{benchmark.py} on the downloaded datasets. Then discuss:
(i)~which backend achieves higher F1/AUCROC per dataset; (ii)~mean and
tail latency of each backend against the 20\,ms SLA; (iii)~whether the
accuracy gap, if any, justifies the latency gap.]

\begin{table*}[!t]
\caption{Detection Accuracy and Steady-State Latency (SLA = 20\,ms).
[TBD: fill from pipeline outputs.]}
\label{tab:steady}
\centering
\begin{tabular}{llccccc}
\hline
Dataset & Backend & F1 & AUCROC & Mean lat. (ms) & p99 lat. (ms) & SLA viol. (\%) \\
\hline
C-MAPSS & LSTM Autoencoder      & [TBD] & [TBD] & [TBD] & [TBD] & [TBD] \\
C-MAPSS & Fast (STL+DMD)        & [TBD] & [TBD] & [TBD] & [TBD] & [TBD] \\
SWaT    & LSTM Autoencoder      & [TBD] & [TBD] & [TBD] & [TBD] & [TBD] \\
SWaT    & Fast (STL+DMD)        & [TBD] & [TBD] & [TBD] & [TBD] & [TBD] \\
\hline
\end{tabular}
\end{table*}

\subsection{Elasticity Under Throughput Spikes}
Table~\ref{tab:stress} reports the stress-protocol results.
[TBD: populate and discuss which backend degrades more gracefully ---
in particular whether the fast backend's lower per-vector compute
translates into a materially lower drop rate at identical queue size and
burst intensity, and how tail latency shifts under saturation.]

\begin{table*}[!t]
\caption{Stress/Elasticity Results Under Unthrottled Burst Ingestion.
[TBD: fill from pipeline outputs.]}
\label{tab:stress}
\centering
\begin{tabular}{llcccc}
\hline
Dataset & Backend & Throughput (vec/s) & Drop rate (\%) & Queue HWM & p99 lat. (ms) \\
\hline
C-MAPSS & LSTM Autoencoder & [TBD] & [TBD] & [TBD] & [TBD] \\
C-MAPSS & Fast (STL+DMD)   & [TBD] & [TBD] & [TBD] & [TBD] \\
SWaT    & LSTM Autoencoder & [TBD] & [TBD] & [TBD] & [TBD] \\
SWaT    & Fast (STL+DMD)   & [TBD] & [TBD] & [TBD] & [TBD] \\
\hline
\end{tabular}
\end{table*}

\subsection{Threats to Validity}
Three limitations qualify the results. First, C-MAPSS anomaly labels are
derived heuristically (final $N$ cycles before failure) rather than being
natively point-labeled as in SWaT; this labeling policy affects absolute
F1 values and cross-dataset comparability. Second, the streaming
architecture is a single-machine multi-threaded simulation rather than a
distributed deployment; absolute latency figures may not transfer to
production hardware without re-benchmarking. Third, the fast backend is a
good-faith reconstruction of the STL+DMD approach described in prior work
rather than a validated reference implementation.

% =========================================================================
\section{Conclusion and Future Work}
% =========================================================================
This paper presented an end-to-end streaming architecture for real-time
industrial MTS anomaly detection under a sub-20\,ms SLA, directly
comparing a deep LSTM autoencoder against a closed-form STL+DMD backend on
identical data and infrastructure. [TBD: one to two sentences stating the
headline empirical conclusion once results are in --- e.g., whether the
fast backend meets the SLA with materially more headroom, and at what
accuracy cost, if any.]

Future work includes: learning an adaptive weighting $\alpha$ in
\eqref{eq:score} rather than fixing it; validating the latency findings on
a distributed ingestion substrate (e.g., Kafka) beyond the single-machine
simulation; and investigating a hybrid mode that ensembles both backends'
scores where the combined cost still fits the SLA budget.

% =========================================================================
\begin{thebibliography}{9}

\bibitem{tao2025realtime}
L.~Tao, M.~Ma, S.~Zhang, and D.~Pei, ``Real-time anomaly detection for
large-scale network devices,'' \emph{IEEE Transactions on Networking},
vol.~33, no.~3, pp. 1326--1337, Jun. 2025.

\bibitem{exarchos2021anomaly}
T.~P. Exarchos \emph{et al.}, ``Anomaly detection in industrial plants
using neural network architectures,'' \emph{IEEE Transactions on
Industrial Informatics}, vol.~17, no.~8, pp. 5432--5441, Aug. 2021.

\bibitem{malhotra2015lstm}
P.~Malhotra, L.~Vig, G.~Shroff, and P.~Agarwal, ``Long short term memory
networks for anomaly detection in time series,'' in \emph{Proc. European
Symposium on Artificial Neural Networks (ESANN)}, 2015, pp. 89--94.

\bibitem{richman2000physiological}
J.~S. Richman and J.~R. Moorman, ``Physiological time-series analysis
using approximate entropy and sample entropy,'' \emph{American Journal of
Physiology-Heart and Circulatory Physiology}, vol. 278, no.~6, pp.
H2039--H2049, Jun. 2000.

\bibitem{saxena2008damage}
A.~Saxena, K.~Goebel, D.~Simon, and N.~Eklund, ``Damage propagation
modeling for aircraft engine run-to-failure simulation,'' in \emph{Proc.
1st International Conference on Prognostics and Health Management (PHM)},
2008, pp. 1--9.

\bibitem{cleveland1990stl}
R.~B. Cleveland, W.~S. Cleveland, J.~E. McRae, and I.~Terpenning, ``STL: A
seasonal-trend decomposition procedure based on loess,'' \emph{Journal of
Official Statistics}, vol.~6, no.~1, pp. 3--73, 1990.

\bibitem{schmid2010dmd}
P.~J. Schmid, ``Dynamic mode decomposition of numerical and experimental
data,'' \emph{Journal of Fluid Mechanics}, vol. 656, pp. 5--28, 2010.

\bibitem{goh2016swat}
J.~Goh, S.~Adepu, K.~N. Junejo, and A.~Mathur, ``A dataset to support
research in the design of secure water treatment systems,'' in \emph{Proc.
International Conference on Critical Information Infrastructures Security
(CRITIS)}, 2016.

% [TBD: add a citation for the VersaGuardian model once a primary source
% is located — see the footnote in Section II-C.]

\end{thebibliography}
